% -----------------------------------------------
% Template for ISMIR LBD Papers
% 2026 version, based on previous ISMIR templates

% Requirements :
% * 2+n page length maximum
% * 10MB maximum file size
% * Copyright note must appear in the bottom left corner of first page
% * Clearer statement about citing own work in anonymized submission
% (see conference website for additional details)
% -----------------------------------------------

\documentclass{article}
\usepackage[T1]{fontenc} % add special characters (e.g., umlaute)
\usepackage[utf8]{inputenc} % set utf-8 as default input encoding
\usepackage[lbd]{ismir}
% !!!PARA VERSAO A SUBMETER USAR LINHA ABAIXO, REMOVER A ACIMA!!!
% \usepackage[lbd,submission]{ismir} 
\usepackage{amsmath,cite,url}
\usepackage{graphicx}
\usepackage{color}
\def\lbd{}	% Flag to use correct LBD settings in the paper, please do not modify this line
\usepackage{listings}
\usepackage{xcolor}

\definecolor{cmtgray}{gray}{0.42}
\lstdefinestyle{layerit}{%
  language=Python,
  basicstyle=\ttfamily\scriptsize,      % 7 pt Courier: 55 characters fit one ISMIR column
  commentstyle=\color{cmtgray}\itshape,
  keywordstyle=\bfseries,
  showstringspaces=false,
  columns=fullflexible, keepspaces=true,
  breaklines=false,                     % every line is short enough that it never has to break
  aboveskip=0pt, belowskip=0pt, xleftmargin=0pt, numbers=none}
\lstset{style=layerit}

% Title. Please use IEEE-compliant title case when specifying the title here,
% as it has implications for the copyright notice
% ------
\title{LayerIt: \linebreak Towards a Framework for Time-Aligned, Composable Music Visualisations}

% Note: Please do NOT use \thanks or a \footnote in any of the author markup

% Single address
% To use with only one author or several with the same address
% ---------------
\oneauthor
  {Fernando Azeredo \hspace{1cm} Ant\'onio S\'a Pinto}
  {Faculdade de Engenharia da Universidade do Porto, Portugal \\
   INESC TEC, Porto, Portugal \\
   {\tt up202308655@fe.up.pt}}

% For the author list in the Creative Common license and PDF metadata, please enter author names.
% Please abbreviate the first names of authors and add 'and' between the second to last and last authors.
\def\authorname{F. Azeredo and A. S\'a Pinto}

% Optional: To use hyperref, uncomment the following.
\usepackage[bookmarks=false,pdfauthor={\authorname},pdfsubject={\papersubject},hidelinks]{hyperref}
% Mind the bookmarks=false option; bookmarks are incompatible with ismir.sty.

\sloppy % please retain sloppy command for improved formatting

\begin{document}

%
\maketitle
%
\begin{abstract}
Music Information Retrieval (MIR) research often requires figures that relate audio representations, symbolic scores, and algorithmic outputs on a common temporal basis. We present LayerIt, a Python framework for generating score-aligned visualisations from independent visual components. It aligns audio, algorithmic data, and symbolic scores on a shared timeline using time-warped scores rendered through \textit{ScoreWarp}, combining them into a structured SVG file. The SVG retains the score's MEI-derived structure and represents each added component as an identifiable group. In a beat-tracking case study, we generate a four-part figure comprising a time-warped score, mel spectrogram, \textit{BeatThis!} beat and downbeat activation (logit) curves, and waveform with the \textit{BeatThis!} output. The example is specified in approximately 60 lines of Python without manual alignment. LayerIt is available as an open-source Python package, providing a reusable basis for score-aligned MIR visualisation workflows.
\end{abstract}
%
\section{Introduction}\label{sec:introduction}

Music can be represented in many ways: as engraved notation, symbolic encodings, or audio recordings of performances. These representations, however, live in different notions of time. Recordings unfold in physical time, measured in seconds or samples; symbolic encodings organize it in abstract units such as beats and measures; and engraved notation renders it as space, in pixels or millimetres. Establishing correspondences between them, a task known as alignment or synchronisation, is central to both music information retrieval (MIR) and digital musicology.

Alignment algorithms, most prominently for the audio-to-score task, must balance accuracy against computational cost and robustness\cite{muller2021SyncToolboxPython}; this alone takes several forms, such as tolerance to performance errors \cite{nakamura2017PerformanceErrorDetection} or to structural mismatches between score and performance \cite{fremerey2010HandlingRepeatsJumps}. Two further axes cut across these choices: where matching takes place, within the audio or symbolic domain, or across the two \cite{mueller2019CrossModalMusicRetrieval}; and the machinery relied on, from classical signal-processing features to learned representations \cite{zeitler2025RobustAccurateAudio}. 

Once computed, an alignment can be put to work in more than one way. One concern is representation: alignments are often stored in ad hoc, single-use formats, and recent modelling relates musical timelines across the graphical, logical, and physical domains so that positions and annotations move between them and the alignment can be reused \cite{hentschel2026TimeAlignModelling}. A separate concern, which this work takes up, is visualization. The prevailing strategy fixes the score's layout and links its elements to instants on a physical timeline through drawn connectors \cite{thomas2012LinkingSheetMusic}, as do interactive tools that track a performance with a moving score cursor \cite{jiang2025PianoPrecisionAdvancing}. Score and physical time thus stay in separate coordinate systems, leaving the reader to bridge them mentally.

An alternative dispenses with anchors altogether: the notation itself is laid out in performance time, each element repositioned according to its performed onset, so that the score stays fully readable as notation while doubling as a visual index into physical time \cite{goebl2025LetsScoreWarpAgain}. This opens the possibility that motivates the present work: a shared performance-time axis onto which further material can be layered, from waveforms and feature curves to other information \cite{pinto2020ShiftIfYou}. Of particular interest are annotations such as beat/downbeat positions or metre: explicit in musical notation, yet recoverable from audio only by estimation.

What is lacking is a simple, programmatic way to compose these layers reproducibly. We present LayerIt: given a score--performance alignment, it places independently generated components onto this shared axis and merges them into a single, self-contained visualization. Each is handled as an independent layer that can be added, combined, or removed, making such composites straightforward to build and adapt. Implemented as an open-source Python library, LayerIt\footnote{\url{https://github.com/FernandoAze/LayerIt}} composes this material where it is already produced — as activation arrays, spectrograms and onset lists — rather than exporting it to a separate drawing tool.

\section{The LayerIt Framework}\label{sec:framework}

% --- P1: what a figure is, and what it costs to make one (~70 w) --------- Opens on the object, not on a disclaimer. "Panel" is introduced here because P2 (panel width) and P4 (per-panel identifiers) both depend on it.
% The boundary --- LayerIt is a renderer, not a synchroniser --- lands in the last sentence, where it reads as a positive contrast rather than an opening negative. Do NOT write "rendered through ScoreWarp": LayerIt never invokes ScoreWarp, it consumes its output. (md, Sec.2 "What's wrong now" 7) Cite roles are deliberate: weigl2020 = the onset/identifier alignment format, goebl2025 = ScoreWarp. (md, Sec.2 P1 Note)
A LayerIt figure is a stack of panels that share one performance-time axis.
Four inputs define it: a performance recording; a score already laid out in performance time by ScoreWarp \cite{goebl2025LetsScoreWarpAgain}, carrying an embedded time axis; a note-level alignment listing performed onsets and the score elements they belong to \cite{weigl2020MultimodalMusicInformation}; and any number of algorithmic outputs.
LayerIt renders that alignment rather than computing it, and emits a single self-contained Scalable Vector Graphics (SVG) document.

% --- P2: the coordinate transformation (~79 w) --------------------------- The section's centre of gravity. This mechanism is currently never stated anywhere in the paper and is the highest-value single addition. (md, Sec.2 P2)
All horizontal placement derives from two numbers read from the warped score's time axis: the span it covers in seconds and its length in pixels.
Their ratio is the only scale factor in the system.
Each layer converts its own native units — samples, frames, or annotated times — into seconds as it loads, with audio read through librosa \cite{mcfee2015LibrosaAudioMusic}, and a shared context maps seconds to pixels identically for all of them.
Each layer controls its own vertical placement; no layer controls its horizontal placement.
% TONE FIX (rev. 3): the earlier version ended "...stays a layer's own business; horizontal placement is never a layer's business at all." Punchy, but too colloquial for the venue. The replacement keeps the parallelism and the design principle without the register slip. Alternatives if this reads flat:
%   "Vertical placement remains layer-specific; horizontal placement is determined once, for all layers."
%   "Vertical placement is defined by each layer; horizontal placement is not."

% TO DECIDE (D2): A / B --- see md, Sec.2 P2 "Accuracy note"
%   A. Keep the four sentences above as they stand. (recommended)
%   B. If words are needed elsewhere, drop sentence 3 and keep the first two plus the last. Costs ~28 w, loses the concreteness that makes the claim checkable by a reader holding the repository.

% --- P3: the interface, and the evidence for it (~73 w) ------------------ THIS PARAGRAPH IS WHAT EARNS "Towards" IN THE TITLE. Do not hedge it, and do not mention the planned shape-based layer taxonomy here or anywhere else in the paper. (md, Sec.0 "the test for any hedge"; Sec.0b) Matplotlib is cited HERE, at first mention, not in P5.
A representation joins a figure as a class implementing two methods: one that loads and validates its data, one that draws it into a Matplotlib \cite{hunter2007Matplotlib2DGraphics} axes, so that existing plotting knowledge transfers unchanged.
A third, optional method lets a layer emit its own SVG group.
Nine layers are implemented on this interface across three families---audio representations, score annotations derived from the alignment, and algorithm outputs---and none of them touches the aligning machinery.

% NOTE (do not delete without reading md, Sec.2 P3 Note 3): the previous draft's phrase "allows additional layer types or rendering backends to be introduced" is deliberately gone. In Matplotlib a "backend" is a swappable renderer behind a stable API; LayerIt has no such seam (to_svg_group returns a hand-assembled SVG string per layer), so the plural names a second backend that does not exist and is disprovable in ninety seconds from the repository.
%
% NOTE (md, Sec.2 P3 Note 2): the Matplotlib sentence above is a FACT, not an analogy --- Layer.draw(self, ax: Axes, ...) puts the axes in the required method's signature and Visualizer.draw calls plt.subplots. Claim only the shared-x-axis half of the subplots comparison, never the automatic-stacking half, unless D5 resolves to (b).

% --- P4: the output, and what it preserves (~78 w) ----------------------- MEI expanded here, at first use in the paper.
The score enters as Music Encoding Initiative (MEI) notation rendered by Verovio \cite{pugin2014VerovioLibraryEngraving} and is not rewritten: its element identifiers and hierarchy survive into the composite.
Each added component is emitted as its own named group, with identifiers rewritten per panel so that repeated layers stay distinguishable, and waveforms, activation curves, onset markers and detected beats remain native SVG primitives---selectable, restylable and animatable after the fact.
%
% TO DECIDE (D1): A / B --- see md, Sec.0 stance clause 1, and Sec.0b
%   A. SHIPS the ~2 h Chromagram.to_svg_group fix. Use the sentence below: it is a statement of design principle, not an admission, and it makes "every layer emits its own named group" universally true. (recommended)
Dense time--frequency representations embed as a raster image inside their group, with their axes drawn over it as SVG.
%   B. Fix does NOT ship. Delete the sentence above and use instead:
% The chromagram and the mel spectrogram are currently raster-only.
%      ("Nine layers are implemented" in P3 stays correct either way, but under B do not claim that every layer emits its own group.)


\section{Demonstration}\label{sec:demonstration}

We demonstrate LayerIt in the figure a beat-tracking analysis calls for: a
recording, a score, and a tracker's frame-level output. The recording and the
tracker's output already lie in physical time; the score does not, and is warped
into it by ScoreWarp from an onset alignment, so its horizontal placement is
computed rather than set by hand. The material is the Prelude from
J.\,S.~Bach's \emph{Prelude and Fugue No.~2 in C minor}, BWV~847, performed by
Glenn Gould; the onsets were aligned to the MEI score in Piano Precision
\cite{jiang2025PianoPrecisionAdvancing}, and the beat and downbeat activations and estimates
come from Beat This!~\cite{foscarin2024BeatThisAccurate} run on the same audio.

\begin{figure}[t]
\begin{lstlisting}
from layerit import Visualizer
from layerit.shapes import Curve, Events, Field
from layerit.sources import (
    MelSpec, Waveform, Onsets, Beats, Downbeats,
    BeatActivation, DownbeatActivation)

fig = Visualizer(audio="Prelude_n2_Cm.wav",
                 score="PreludeN2_WScore.svg",
                 maps="OnsetsMei_PreludioN2.maps.json",
                 beats="beat_Bach.npz")

marks = Events(Onsets(), color="black", line_width=0.3)

fig.add_panel(Field(MelSpec(freq_window=(100, 1500)),
                    color_map="summer"), marks)
fig.add_panel(Curve(BeatActivation(), line_width=0.7),
              Curve(DownbeatActivation(),
                    line_width=0.7),
              marks)
fig.add_panel(Curve(Waveform(), color="magenta"),
              Events(Beats(), color="red"),
              Events(Downbeats(), color="blue"), marks)

fig.compose("LDB_FIG.svg")
\end{lstlisting}
\caption{WIP: The complete specification of Figure~\ref{fig:composite}.
The four inputs are declared once; each \texttt{add\_panel} call combines the layers of one panel; \texttt{compose} stacks them on the axis read from the warped score.
A layer is one of four shapes --- a field, a curve, events or intervals --- holding the data it draws, so drawing is written once per shape rather than once per representation, and a new representation is added by choosing the shape that fits it.
Panel order is the only placement the script states.}
\label{fig:script}
\end{figure}

Figure~\ref{fig:script} shows the script that produces Figure~\ref{fig:composite}: four components, each declared in five to seven lines against the same interface, plus the composition step that stacks them---about sixty lines in total.
The horizontal axis is established once, from the warped score; the panel order is the only placement the script states explicitly.

Substituting a different beat tracker changes one path in the script; a different performance additionally requires its own alignment and warped score, which LayerIt does not compute.

\section{Conclusion}\label{sec:conclusion}

LayerIt provides a programmatic workflow for generating score-aligned MIR figures from independently specified components. Its structured SVG output preserves the score's semantic information and the identity of added components, rather than flattening the figure into a single raster image. The current implementation is open source and provides a stable core for additional layer types. Planned extensions include tempograms and spectrogram derivatives, time-alignment backends beyond \textit{ScoreWarp}, and annotation export. Integrating an upstream synchronisation toolkit such as \textit{Sync Toolbox}, and supporting interchange data such as match files, would broaden the set of alignment inputs available to the framework. The structured SVG representation may also support interactive, playback-synchronised viewers. By making the figure specification re-runnable, LayerIt supports transparent and reusable visualisation workflows.


%\newpage

%\section{AI Usage Statement}

%Large language models (LLMs) were used in this work as a writing and code development aid. Specifically: (1)~LLMs were used to help structure and write this LBD paper, including organization, phrasing, and clarity improvements; (2)~LLMs were used to assist in the development of the LayerIt codebase. The core research methodology, conceptual contributions, framework architecture, and scientific validation remain the work of the authors. All code, design decisions, and experimental results have been authored and verified by the authors themselves.

% For bibtex users:
\bibliography{ISMIRtemplate_ASP}

\end{document}
